\begin{enquestion}{20}{
    An artillery reconnaissance unit is analyzing the impact dispersion of unguided projectiles on a test range.
    Based on wind tunnel models and telemetry data, the joint probability density function (PDF) for the longitudinal and lateral miss distances ($X$, $Y$)
    relative to the target center (in tens of meters) is given by:
    \begin{align*}
        f_{X,Y}(x,y) =  \begin{cases}
                            C \cdot (x + 4y), & \text{if } 0 \le x \le 4 \text{ and } 0 \le y \le 1 \\
                            0, & \text{otherwise,}
                        \end{cases}
    \end{align*}
    where $C$ is a constant.
}

    \subquestion{6}{Determine the value of $C$ so that $f_{X,Y}(x,y)$ satisfies the properties of a valid joint probability density function.}
    \ensolution{
        The function $f_{X,Y}(x,y)$ is a valid joint probability density function if the following two conditions hold:
        \begin{enumerate}
            \item $f_{X,Y}(x,y) \ge 0$ for all possible values of $x$ and $y$, and
            \item The volume bounded by $f$ and the $xy$-plane is equal to $1$:
            \[
                \int_{-\infty}^{\infty} \int_{-\infty}^{\infty} f_{X,Y}(x,y)\dx\dy = 1.\rubric{1}
            \]
        \end{enumerate}
        Since $x + 4y \ge 0$ for $0 \le x \le 4$ and $0 \le y \le 1$, the first condition holds if and only if $C \ge 0$. \rubric{1}
        
        The second condition holds if
        \begin{align*}
            1   &= \int_{-\infty}^{\infty} \int_{-\infty}^{\infty} f_{X,Y}(x,y)\dx\dy \\
                &= \int_0^1 \int_0^4 C \cdot (x + 4y)\dx\dy \\
                &= C \cdot \int_0^1 \left[\int_0^4 (x + 4y)\dx\right] \dy \\
                &= C \cdot \int_0^1 \left[\frac{1}{2}x^2 + 4xy\right]_0^4 \dy \\
                &= C \cdot \int_0^1 (8 + 16y) \dy \\
                &= C \cdot \left[8y + 8y^2\right]_0^1 \\
                &= C \cdot (8 + 8) = 16 \cdot C. \rubric{3}
        \end{align*}
        Thus, the constant $C$ must be equal to $\frac{1}{16}$ in order for $f_{X,Y}(x,y)$ to be a valid joint probability density function. \rubric{1}
    }
    
    \subquestion{6}{Compute the marginal PDFs $f_X(x)$ and $f_Y(y)$, for $X$ and $Y$ respectively.}
    \ensolution{
        \begin{itemize}
            \item For $0 \le x \le 4$, the marginal PDF of $X$ is given by 
            \begin{align*}
                f_{X}(x) = \int_{-\infty}^{\infty} f_{X,Y}(x,y) \dy &= \int_0^1 \frac{1}{16} \cdot (x + 4y) \dy \\
                                                                   &= \frac{1}{16} \left[xy + 2y^2\right]_{0}^1 \\
                                                                   &= \frac{1}{16}x + \frac{1}{8}, \rubric{3}
            \end{align*}
            and $f_X(x) = 0$ otherwise.
            
            \item For $0 \le y \le 1$, the marginal PDF of $Y$ is given by 
            \begin{align*}
                f_{Y}(y) = \int_{-\infty}^{\infty} f_{X,Y}(x,y) \dx &= \int_0^4 \frac{1}{16} \cdot (x + 4y) \dx \\
                                                                   &= \frac{1}{16} \left[\frac{1}{2}x^2 + 4xy\right]_{0}^4 \\
                                                                   &= \frac{1}{16}(8 + 16y) = y + \frac{1}{2}, \rubric{3}
            \end{align*}
            and $f_Y(y) = 0$ otherwise.
        \end{itemize}
    }

    \subquestion{8}{Are $X$ and $Y$ independent random variables? Justify your answer with a computation.}
    \ensolution{
        \textbf{Method 1 (Using the definition of independence):} \\
        Two continuous random variables $X$ and $Y$ are independent if and only if $f_{X,Y}(x,y) = f_X(x) \cdot f_Y(y)$ for all $(x,y) \in \mathbb{R}^2$. \rubric{2}

        For $0 \le x \le 4$ and $0 \le y \le 1$, the product of the marginal PDFs is:
        \begin{align*}
            f_X(x) \cdot f_Y(y) &= \left(\frac{1}{16}x + \frac{1}{8}\right)\left(y + \frac{1}{2}\right) \\
                                &= \frac{1}{16}xy + \frac{1}{32}x + \frac{1}{8}y + \frac{1}{16}. \rubric{4}
        \end{align*}
        Comparing this to the joint PDF $f_{X,Y}(x,y) = \frac{1}{16}x + \frac{1}{4}y$, we see that $f_{X,Y}(x,y) \neq f_X(x) \cdot f_Y(y)$ (for example, at $(x,y) = (0,0)$, $f_{X,Y}(0,0) = 0$ whereas $f_X(0)f_Y(0) = \frac{1}{16}$). \rubric{1}
        
        Hence, $X$ and $Y$ are \emph{not} independent random variables. \rubric{1}

        \vspace{0.5em}
        \textbf{Method 2 (Alternative via covariance):} \\
        If $X$ and $Y$ were independent, their covariance $\Cov{X}{Y} = E(XY) - E(X)\cdot E(Y)$ would have to be equal to zero. \rubric{1}
        
        Computing the required expectations:
        \begin{align*}
            E(X) &= \int_{0}^{4} x \cdot \left(\frac{1}{16}x + \frac{1}{8}\right) \dx = \left[\frac{1}{48}x^3 + \frac{1}{16}x^2\right]_{0}^{4} = \frac{4}{3} + 1 = \frac{7}{3}. \rubric{2} \\
            E(Y) &= \int_{0}^{1} y \cdot \left(y + \frac{1}{2}\right) \dy = \left[\frac{1}{3}y^3 + \frac{1}{4}y^2\right]_{0}^{1} = \frac{1}{3} + \frac{1}{4} = \frac{7}{12}. \rubric{2} \\
            E(XY) &= \int_{0}^{1} \int_{0}^{4} \frac{1}{16}(x^2y + 4xy^2)\dx\dy = \frac{1}{16} \int_0^1 \left[\frac{1}{3}x^3y + 2x^2y^2\right]_{0}^{4} \dy \\
                  &= \int_0^1 \left(\frac{4}{3}y + 2y^2\right) \dy = \left[\frac{2}{3}y^2 + \frac{2}{3}y^3\right]_0^1 = \frac{4}{3}. \rubric{2}
        \end{align*}
        Therefore,
        \[
            \Cov{X}{Y} = \frac{4}{3} - \frac{7}{3} \cdot \frac{7}{12} = \frac{48 - 49}{36} = -\frac{1}{36} \neq 0.
        \]
        Since $\Cov{X}{Y} \neq 0$, $X$ and $Y$ cannot be independent random variables. \rubric{1}
    }
\end{enquestion}